% Transcription — fonds Alexandre Grothendieck, shelfmark 162-6, batch 1 (pages 1-5).
% Established from the facsimile archives/batches/162-6/batch-01.pdf (a mirror of
% https://grothendieck.umontpellier.fr/162-6.pdf), per the skill
% transcribe-grothendieck.
% Five small cards in pencil, written landscape, one subject throughout: the
% archimedean gamma factor of an L-function, its recipe from the Hodge numbers,
% the real Weil group, and the group S_E through which the algebraic Hecke
% characters of E factor. The inventory calls the folder « Documents isolés »
% and describes nothing; the five cards in fact hang together and are
% transcribed here in the archivists' order, which is also a readable one.
% Pass: Opus 5 (claude-opus-5), 2026-08-16 — first pass, unchecked against the pages by a human.
\documentclass[11pt,a4paper]{article}
% Le préambule commun à toutes les transcriptions du fonds Grothendieck.
%
% Il tient en une idée : **l'incertitude se compose, elle ne se résout pas.**
% Une transcription qui lisse un mot douteux a détruit la seule chose pour
% laquelle on la faisait — un lecteur doit pouvoir savoir, à chaque endroit,
% ce qui était sur le feuillet et ce qui a été ajouté. Les six macros qui
% suivent sont donc l'appareil critique, et non de la décoration.
%
% Les mêmes commandes sont reconnues par `scripts/render.mjs`, qui produit la
% vue de lecture du site. Une macro ajoutée ici doit l'être aussi là-bas,
% faute de quoi le rendu échouera bruyamment — ce qui est voulu : un rendu
% silencieusement faux est pire qu'un rendu absent.

\ProvidesPackage{grothendieck}[2026/08/08 Transcriptions du fonds Grothendieck]

\usepackage{amsmath,amssymb,amsthm}
\usepackage{mathrsfs}
\usepackage{stmaryrd}
% Les diagrammes commutatifs sont partout dans ce fonds : c'est la langue
% même de la géométrie algébrique de Grothendieck, et une transcription qui
% les rendrait en prose ne transcrirait rien.
\usepackage{tikz-cd}
% Français par défaut — c'est la langue des feuillets — mais l'anglais est
% chargé aussi : la traduction et le résumé sont des documents anglais, et la
% typographie française (espaces avant les deux-points) y serait une faute.
% Ces documents basculent par \selectlanguage{english} en tête de corps.
\usepackage[english,french]{babel}
\usepackage{fontspec}
\usepackage{geometry}
\geometry{a4paper,margin=25mm,marginparwidth=28mm}
\usepackage{marginnote}
\usepackage{xcolor}
% ulem, not soul. `soul`'s \st and \ul are built on a character-by-character
% scan that XeLaTeX's font machinery breaks: the document fails with an
% undefined control sequence rather than degrading. ulem does the same two jobs
% and is Unicode-engine safe. `normalem` stops it hijacking \emph, which the
% transcriptions use for Grothendieck's own emphasis.
\usepackage[normalem]{ulem}
\usepackage{hyperref}
\hypersetup{colorlinks=true,linkcolor=black,urlcolor=black,pdfborder={0 0 0}}

% --- Métadonnées du lot -----------------------------------------------------
% Reprises telles quelles par le script de rendu, qui les affiche en tête de
% la vue de lecture. Elles fixent ce que le fichier prétend couvrir : sans
% elles, une transcription flotte sans point d'attache dans un fonds de seize
% mille feuillets.

\newcommand{\folder}[1]{\gdef\grothFolder{#1}}
\newcommand{\batch}[1]{\gdef\grothBatch{#1}}
\newcommand{\leaves}[2]{\gdef\grothFirst{#1}\gdef\grothLast{#2}}
\newcommand{\foldertitle}[1]{\gdef\grothFoldertitle{#1}}
\gdef\grothFolder{?}\gdef\grothBatch{?}\gdef\grothFirst{?}\gdef\grothLast{?}\gdef\grothFoldertitle{}

% --- Le résumé --------------------------------------------------------------
%
% Il ouvre la lecture modernisée, sous le titre « Résumé », et il est composé
% en petit. Ce n'est pas un ornement : c'est le seul passage du document qui
% ne soit pas des mathématiques, et un lecteur doit pouvoir le reconnaître
% avant de l'avoir lu — puis le sauter s'il connaît déjà le sujet, ou s'y
% arrêter s'il ne le connaît pas. Le filet qui le suit marque la reprise du
% corps.

\newenvironment{resume}
  {\small\setlength{\parskip}{0.45em}}
  {\par\medskip\hrule\medskip}

% --- Mots-clés ---------------------------------------------------------------
%
% En anglais, en clôture du résumé de la lecture modernisée : le vocabulaire
% moderne sous lequel chercher ce que les feuillets construisent. Le manifeste
% du site (`npm run manifest`) les extrait pour étiqueter la cote entière —
% cette ligne est la source unique des tags, il n'y a pas de fichier à côté.
\newcommand{\keywords}[1]{\par\smallskip\noindent
  {\footnotesize\textsc{Keywords} --- \textit{#1}}\par}

% --- L'appareil critique ----------------------------------------------------

% Le numéro de feuillet, en marge. C'est l'ancre qui permet de régler un
% désaccord sur une formule en regardant *un* feuillet plutôt qu'une chemise
% de six cents. Le site s'en sert aussi pour tourner les pages du fac-similé
% quand on fait défiler la transcription.
\newcommand{\leaf}[1]{%
  \marginnote{\footnotesize\color{gray!70}\textsf{#1}}%
  \label{leaf:#1}\ignorespaces}

% Illisible : on ne devine pas. Un mot inventé qui se lit comme les autres est
% le pire résultat possible.
\newcommand{\ill}{\textcolor{red!60!black}{[\,\ldots\,]}}

% Lecture douteuse : le mot est donné, et signalé comme incertain.
\newcommand{\uncertain}[1]{\uline{#1}}

% Ajout de l'éditeur — une lettre manquante, un mot sous-entendu.
\newcommand{\add}[1]{\textcolor{blue!50!black}{[#1]}}

% Biffé par Grothendieck lui-même. Ce qu'il a rayé fait partie du document :
% une rature dit souvent où la pensée a changé de direction.
\newcommand{\struck}[1]{\sout{#1}}

% Note du transcripteur, sur l'état matériel du feuillet.
\newcommand{\note}[1]{\par\smallskip{\small\color{gray!60!black}\textit{[#1]}}\par\smallskip}

% Note marginale de Grothendieck, distincte du corps du texte.
\newcommand{\marginal}[1]{\par\smallskip\noindent\hspace{1em}%
  {\small\color{gray!60!black}#1}\par\smallskip}

% --- Titre ------------------------------------------------------------------

\AtBeginDocument{%
  \begin{center}
    {\large\bfseries Fonds Alexandre Grothendieck --- cote n\textsuperscript{o}~\grothFolder}\\[2pt]
    {\itshape\grothFoldertitle}\\[4pt]
    {\small feuillets \grothFirst--\grothLast\ (lot \grothBatch)}\\[2pt]
    {\footnotesize\color{gray!60!black}Transcription établie d'après le fac-similé numérisé
     par l'Université de Montpellier.\\
     Les lectures incertaines sont soulignées, les illisibles notées [\,\ldots\,],
     les ajouts entre crochets.}
  \end{center}
  \vspace{1em}}

\endinput


\folder{162-6}
\batch{1}
\pages{1}{5}
\dating{s.d. — le groupe « Dossiers rassemblés par Grothendieck sur différentes thématiques » (161-1 à 162-6) est daté [après 1961-vers 1977]}
\watermark{Édition de démonstration}
\foldertitle{[Documents isolés] : notes manuscrites (s.d.)}

\begin{document}

\page{1}

\[
\gamma_{\mathbb{R}}(s) = \pi^{\frac{s}{2}}\,\Gamma\!\left(\frac{s}{2}\right)
\]
\[
\gamma_{\mathbb{C}}(s) = \tfrac{1}{2}\,\ill\,(2\pi)^{s}\,\Gamma(s)
\]

\note{les deux exposants sont écrits sans signe : la page porte
$\pi^{s/2}$ et $(2\pi)^{s}$, non $\pi^{-s/2}$ et $(2\pi)^{-s}$. On les
transcrit tels quels. Entre le $\tfrac{1}{2}$ et la parenthèse un symbole a été
tracé puis barré, illisible}

\emph{(un trait horizontal sépare ce qui précède de ce qui suit)}

Cas $\mathbb{R}$, avec $n = 2p$ :
\[
\prod_{p < q} \gamma_{\mathbb{C}}(s-p)^{h(p,q)} \cdot
\begin{cases}
h_{p}^{+} : F = (-1)^{p}\\
h_{p}^{-} : F = -(-1)^{p}\\
\gamma_{\mathbb{R}}(s-p)^{h_{p}^{+}}\ \gamma_{\mathbb{R}}(s-p+1)^{h_{p}^{-}}
\end{cases}
\]

\note{$F$ est l'involution qui agit sur la partie de poids $p=q$ ; les deux
lignes qui la concernent donnent la règle de partage de $h^{p,p}$ en
$h_{p}^{+}$ et $h_{p}^{-}$, et la troisième le facteur que chacune contribue.
Dans la seconde, le signe devant $(-1)^{p}$ est surchargé et un signe a été
entouré puis abandonné}

Cas $\mathbb{C}$ :
\[
\prod_{p,q} \gamma_{\mathbb{C}}\bigl(s - \inf(p,q)\bigr)^{h(p,q)}
\]

\page{2}

\[
\mathrm{gal}(E/\mathbb{Q}) \;-\; C(E)
\]

\note{le $E$ et le $\mathbb{Q}$ sont chacun repassés ; le trait entre les deux
membres n'est pas une flèche}

$W_{\mathbb{R}}$ :
\[
\begin{array}{c|l}
\sigma & \sigma z \sigma^{-1} = \bar{z}\\
\mathbb{C}^{*} & \sigma^{2} = (-1)
\end{array}
\]
avec, sous les deux termes de gauche, l'accolade $\mathbb{Z}/2$, $\mathbb{C}^{*}$.

\note{le $-1$ est entouré}

\[
H^{2}(\mathbb{Z}/2, \mathbb{C}^{*}) = \mathbb{Z}/2
\]

\emph{(trait horizontal)}

$M$, puis $R - \mathbb{C}$, en regard de quoi : \struck{$H(E)$} suivi en
exposant de $(p,q)$, puis

\[
H(D(E), \mathbb{R}) \quad \sigma \quad \Big|\ \sigma\ \Big|\ \sigma^{2} = +1
\]

et, sous le $\sigma$ du milieu, un $S$ ; à gauche de celui-ci un terme barré et
illisible : \struck{\ill}

\note{la dernière ligne de la carte est coupée par le bord de la feuille : on y
lit encore un $c_{\mathbb{R}}$, un signe de produit tensoriel et un signe
d'égalité, le reste manque — \ill}

\page{3}

\[
N : S_{E} \longrightarrow S_{\mathbb{Q}} \longrightarrow \mathbb{G}_{m}
\]
\[
(x, (y_{\alpha})) \longmapsto x \cdot \prod_{\alpha \neq \infty} \lVert y_{\alpha} \rVert \cdot \prod_{\alpha = \infty} \mathrm{sign}(y_{\alpha})
\]

\[
C(E) \longrightarrow S_{E}(\mathbb{A}) \longrightarrow S_{E}(\mathbb{R}) \longrightarrow \mathbb{R}^{*}
\]

\note{une flèche courbe joint $C(E)$ au dernier terme et porte, en dessous et
souligné deux fois, $\lVert x \rVert$ : c'est la composée qui est nommée}

\emph{(trait horizontal, puis, seul en bas de la carte :)}
\[
(-1 \neq \lVert -1 \rVert)
\]

\page{4}

\[
L_{v}(M \otimes T, s) = L_{v}(M)(s \pm 1)
\]
\[
\gamma_{\mathbb{R}}(s)\, \gamma_{\mathbb{R}}(s+1) = \gamma_{\mathbb{C}}(s)
\]

\note{les deux lignes sont écrites en biais, montantes, et rien d'autre ne
figure sur la carte}

\page{5}

En haut à gauche, encadré : $E/\mathbb{Q}$.

\begin{tikzcd}[column sep=small, row sep=small, nodes={font=\scriptsize}]
& & & \mathbb{G}_{m}\\
E^{*} \arrow[r] & \left(\prod_{E/\mathbb{Q}} \mathbb{G}_{m}\right)\big/U \times I(E)/I^{\circ}(E) \arrow[rr] & & S_{E} \arrow[u]\\
E^{*} \arrow[r] \arrow[u, no head] & I(E) \times I(E)/I^{\circ}(E) \arrow[rr] & & S_{E}(\mathbb{A})\\
E^{*} \arrow[rr] \arrow[u, no head] & & I(E) \arrow[r] & C(E) \arrow[u]
\end{tikzcd}

\note{les deux traits verticaux de la colonne de gauche sont les doubles barres
d'égalité de la page, rendues ici par des liaisons sans tête : c'est le même
$E^{*}$ aux trois lignes. Entre la deuxième et la troisième ligne, au milieu,
deux flèches obliques montantes portent les chiffres 2 (vers la gauche) et 1
(vers la droite) ; leur origine n'est pas fixée sur la carte et elles ne sont
pas reportées dans le diagramme. À droite, une accolade réunit la colonne, et
$S_{E}(\mathbb{Q})$, écrit au-dessus de $S_{E}(\mathbb{A})$, est barré}

\emph{(sous un trait horizontal, en bas de la carte :)}
\struck{$S_{E} \longrightarrow \mathbb{G}$}, puis
$S_{E} \longrightarrow \mathbb{G}_{m}$, puis \ill

\note{la fin de cette dernière ligne sort de la carte : on y distingue un $h$ et
une parenthèse ouvrante suivie d'un $E$, rien de plus}

\end{document}
